\section{话题、群体与干预部署}
\label{app:impl}

\paragraph{话题。}
主要研究及剂量反应和规模分析中使用的四个话题通过简短的政策描述和锚定在$-5$（强烈反对）和$+5$（强烈支持）的立场量表进行操作化。\emph{疫苗强制}：政府是否应为公共部门就业和学校入学强制接种疫苗。\emph{气候政策}：政府是否应在经济成本下积极推进绿色能源转型和碳定价。\emph{移民}：国家是否应采取更开放的移民政策，扩大合法途径和难民接纳。\emph{转基因食品安全}：转基因食品是否应在现有监管下广泛采用，还是应作为企业控制的风险予以拒绝。完整话题描述随配置一同发布。

\paragraph{群体特征。}
代表性群体将每个信念维度从$\mathcal{N}(5.5, 2.0)$采样，截断到$[1, 10]$。目标群体每次修改单一维度：\emph{低信任}设$D_1 \sim \mathcal{N}(2.5, 1.0)$，\emph{高信任}设$D_1 \sim \mathcal{N}(8.5, 1.0)$，\emph{高共情}设$D_2 \sim \mathcal{N}(8.5, 1.0)$，\emph{权威主义}设$D_3 \sim \mathcal{N}(8.5, 1.0)$，\emph{高分析}设$D_5 \sim \mathcal{N}(8.5, 1.0)$，\emph{低分析}设$D_5 \sim \mathcal{N}(2.5, 1.0)$。所有其他维度保持代表性分布。初始立场从$\mathcal{N}(1.5, 1.5)$采样并四舍五入为$[-5, +5]$中的整数，模拟讨论开始时温和支持主流政策方向的群体。

\paragraph{干预部署。}
预防性干预（预防接种、元认知反思）通过在模拟开始时向智能体系统提示附加额外的信念描述来部署；附加文本与智能体本身共享，使智能体\emph{作为}具有该倾向的人进行推理，而非被外部指示。反应性干预（事实核查、延迟预防接种、社区举报、情绪透明）在配置的轮次窗口期间作为标记的上下文块注入用户提示中。
